\documentclass[11pt]{article}

\usepackage[margin=2.5cm]{geometry}
\usepackage{amsmath}
\usepackage{booktabs}
\usepackage{siunitx}
\usepackage{xcolor}
\usepackage[hidelinks]{hyperref}

\newcommand{\dg}{^{\circ}}

\title{Stage 1 --- Minimum-Jerk Motion Generation}
\author{Vision-Based Optical Simulation}
\date{}

\begin{document}
\maketitle

\noindent
This document describes how the Stage-1 elbow-curl \emph{motion} is generated:
a minimum-jerk baseline with a progressive, physiologically motivated fatigue
perturbation.

%==============================================================
\section{Scientific Goal}

The goal is to generate physiologically plausible repetitive elbow-curl motion
\emph{without} motion capture. Instead of recording a person, the elbow angle is
generated from a motor-control law: \textbf{minimum jerk}. This gives a clean,
human-like baseline. Controlled fatigue effects are then added so the motion
becomes slower, smaller in range, shakier, and more asymmetric over repetitions.

\[
\text{minimum-jerk law} \;\rightarrow\; \text{clean elbow curl}
\;\rightarrow\; \text{progressive fatigue perturbation}.
\]

\paragraph{Provenance.} \textit{Inspired by} Flash \& Hogan (1985) and Hogan
(1984) --- the choice of minimum jerk as the model of natural, unconstrained
human movement is theirs. The idea of layering progressive fatigue degradation
on that baseline is \textit{grounded in} the repetitive-task fatigue literature
(Potvin 1997; Savin et al.\ 2021); combining the two into a synthetic generator
is this project's own design.

%==============================================================
\section{Minimum-Jerk Motion}

For one movement from $q_{\text{start}}$ to $q_{\text{end}}$ over duration $T$,
the normalized time is
\[
s = \frac{\tau}{T}, \qquad s \in [0,1],
\]
and the elbow angle follows the fifth-order minimum-jerk polynomial
\begin{equation}
q(\tau) = q_{\text{start}} + \bigl(q_{\text{end}} - q_{\text{start}}\bigr)
\left( 10\,s^{3} - 15\,s^{4} + 6\,s^{5} \right).
\end{equation}

This polynomial is used because it guarantees the properties in
Table~\ref{tab:mj-props}.

\begin{table}[h]
\centering
\begin{tabular}{ll}
\toprule
\textbf{Property} & \textbf{Meaning} \\
\midrule
zero velocity at endpoints & no sudden start/stop \\
zero acceleration at endpoints & no artificial acceleration spike \\
single bell-shaped velocity profile & close to observed human reaching/curl motion \\
smooth position curve & suitable baseline for OpenSim kinematics \\
\bottomrule
\end{tabular}
\caption{Properties guaranteed by the fifth-order minimum-jerk polynomial.}
\label{tab:mj-props}
\end{table}

The analytic derivatives are
\begin{align}
\dot q(\tau)  &= \bigl(q_{\text{end}} - q_{\text{start}}\bigr)
\,\frac{30\,s^{2} - 60\,s^{3} + 30\,s^{4}}{T}, \\[4pt]
\ddot q(\tau) &= \bigl(q_{\text{end}} - q_{\text{start}}\bigr)
\,\frac{60\,s - 180\,s^{2} + 120\,s^{3}}{T^{2}}.
\end{align}

\paragraph{Consistency note.} After tremor and drift are added (Section~\ref{sec:fatigue}),
these analytic derivatives are no longer the derivatives of the final written
signal. The generator therefore computes velocity and acceleration by finite
differences of the final elbow angle that is actually written to the motion file,
so the kinematics remain self-consistent.

\paragraph{Provenance.} \textit{Directly from} Flash \& Hogan (1985): the
fifth-order polynomial, its zero velocity/acceleration at the endpoints, and the
bell-shaped velocity profile are their result; Hogan (1984) applies it to
single-joint (elbow) movement. The analytic derivatives follow by
differentiation. The finite-difference \emph{consistency} step (recomputing
$\dot q,\ddot q$ from the written signal after perturbation) is this project's
implementation choice, not from a paper.

%==============================================================
\section{One Repetition}

One elbow-curl repetition is built from four segments:
\begin{center}
\begin{tabular}{ll}
\toprule
\textbf{Segment} & \textbf{Description} \\
\midrule
$20\dg \rightarrow 120\dg$ & flexion / concentric phase \\
$120\dg$ hold & top pause \\
$120\dg \rightarrow 20\dg$ & extension / eccentric phase \\
$20\dg$ hold & bottom pause \\
\bottomrule
\end{tabular}
\end{center}

The default canonical settings are listed in Table~\ref{tab:params}. The
$20\dg$--$120\dg$ range avoids full extension and gives a realistic
controlled curl range for the simplified arm26 model.

\begin{table}[h]
\centering
\begin{tabular}{ll}
\toprule
\textbf{Parameter} & \textbf{Default} \\
\midrule
sampling rate & \SI{100}{\hertz} \\
start angle $q_0$ & $20\dg$ \\
fresh target $q_f$ & $120\dg$ \\
shoulder angle & $20\dg$ (fixed) \\
flexion duration & \SI{1.5}{\second} \\
extension duration & \SI{1.5}{\second} \\
top / bottom pause & \SI{0.2}{\second} \\
\bottomrule
\end{tabular}
\caption{Canonical motion parameters.}
\label{tab:params}
\end{table}

\paragraph{Provenance.} \textit{Project implementation.} The four-segment
concentric/eccentric structure with top/bottom pauses and the chosen timings are
our design. The $20\dg$--$120\dg$ range is \textit{informed by} controlled-curl
ROM in the arm26 literature (e.g.\ Attarieh et al.\ 2025 used $10\dg$--$140\dg$)
and chosen to avoid full extension; the exact values are our setting.

%==============================================================
\section{Fatigue Model}
\label{sec:fatigue}

Fatigue is represented by a normalized level $f \in [0,1]$. The generator uses a
\textbf{saturating} fatigue curve (not a straight line),
\begin{equation}
f(x) = \frac{1 - e^{-r x}}{1 - e^{-r}},
\end{equation}
where $x$ is repetition progress from $0$ to $1$ and the fatigue rate is
$r = 1.2$. This makes the decline less front-loaded over short protocols while
still allowing faster or slower fatigue.

The fatigue level $f$ drives five effects (Table~\ref{tab:fatigue}).

\begin{table}[h]
\centering
\begin{tabular}{lll}
\toprule
\textbf{Effect} & \textbf{Parameter} & \textbf{What happens} \\
\midrule
slower movement & \texttt{dur\_gain} & movement duration increases up to \SI{40}{\percent} \\
reduced range of motion & \texttt{rom\_loss} & target angle drops from $120\dg$ toward $100\dg$ \\
timing asymmetry & \texttt{asym} & flexion slows, extension relatively faster \\
tremor / less smoothness & \texttt{tremor\_amp} & \SIrange{6}{12}{\hertz} band-limited tremor added \\
endpoint drift & \texttt{drift\_amp} & slow \SIrange{0.1}{0.5}{\hertz} turnaround variability \\
\bottomrule
\end{tabular}
\caption{Fatigue effects controlled by $f$.}
\label{tab:fatigue}
\end{table}

\noindent
The tremor and drift perturbations are \textbf{band-limited sums of sinusoids},
not white noise. This matters because white noise would create unrealistic jerk
and unstable finite-difference derivatives. The tremor occupies a physiological
$\SIrange{6}{12}{\hertz}$ band; the slow drift occupies $\SIrange{0.1}{0.5}{\hertz}$
and produces realistic turnaround variability.

\paragraph{Optional motion variants.}
Two settings shape the high-frequency content without changing the underlying
curl: \texttt{clip\_to\_rom} bounds the final angle to $[q_0, q_f]$, and
\texttt{lowpass\_hz} applies a low-pass filter to produce a lower-noise,
inverse-dynamics-safe motion. The default keeps the full physiological
variability; the velocity and acceleration are always recomputed from the final
written angle, including any clipping or filtering.

\paragraph{Provenance.} The \emph{saturating} fatigue curve is \textit{inspired
by} the three-compartment fatigue model of Xia \& Frey-Law (2008) and Frey-Law
et al.\ (2012) (accumulation toward a plateau). The five degradation effects are
\textit{taken from} the repetitive-task fatigue literature: slowing and reduced
peak velocity / ROM (Potvin 1997; Potvin \& Bent), timing asymmetry and
movement-variability / endpoint drift (Savin et al.\ 2021; Yang et al.\ 2018;
Sheikhhoseini et al.\ 2025). The \emph{band-limited} (sum-of-sinusoids, not white)
implementation of tremor/drift and the \texttt{clip\_to\_rom} / \texttt{lowpass\_hz}
options are this project's implementation choices, not from a specific paper.

%==============================================================
\section{Quality-Control Validation}

Each generated motion is validated against minimum-jerk theory and broad
literature expectations. The checks are:
\begin{itemize}\itemsep2pt
  \item fresh peak velocity close to the minimum-jerk prediction
        $\dot q_{\max} = 1.875\,(q_f - q_0)/T$;
  \item fresh velocity profile approximately symmetric (time-to-peak $\approx \SI{50}{\percent}$);
  \item fresh range of motion matches the commanded ROM;
  \item fatigue causes peak velocity to decrease over repetitions;
  \item fatigue causes range of motion to decrease over repetitions;
  \item fatigue causes movement duration to increase over repetitions;
  \item peak elbow velocity remains in a plausible human range ($<300\,\dg\,\mathrm{s}^{-1}$).
\end{itemize}

The fatigue trends are tested as least-squares regression \emph{slopes} across all
repetitions, not by comparing only the first and last repetition. This is more
robust because endpoint drift is deliberate and can mask a small difference in a
two-point comparison.

\paragraph{Provenance.} The peak-to-average velocity ratio $1.875$ is the
minimum-jerk value confirmed empirically by Flash \& Hogan (1985). The
fatigue-trend expectations (peak velocity and ROM decrease, duration increases)
are \textit{from} the fatigue literature (Potvin 1997; Potvin \& Bent; Savin et
al.\ 2021). The regression-slope test procedure and the numeric thresholds are
this project's implementation.

%==============================================================
\section{What Changed From the Original Motion}

\begin{table}[h]
\centering
\begin{tabular}{lll}
\toprule
\textbf{Feature} & \textbf{Original motion} & \textbf{Improved generator} \\
\midrule
elbow ROM   & $0\dg \rightarrow 80\dg$ & $20\dg \rightarrow 120\dg$, declining with fatigue \\
shoulder    & $0\dg$ & $20\dg$ fixed stabilization \\
movement model & minimum jerk only & minimum jerk $+$ fatigue perturbation \\
fatigue     & none & slower, smaller ROM, asymmetric, tremor, drift \\
noise       & none & band-limited tremor/drift (not white noise) \\
derivatives & not provided & finite differences of the written signal \\
validation  & manual inspection & automatic QC report \\
\bottomrule
\end{tabular}
\caption{Changes from the original motion to the improved minimum-jerk generator.}
\end{table}

\paragraph{Provenance.} \textit{Project engineering} --- this is a comparison of
the improved generator against the original motion file produced earlier in the
project, not content drawn from a paper.

%==============================================================
\section*{References}
\begin{description}\itemsep2pt
\item[Hogan (1984).] An organizing principle for a class of voluntary movements.
  \emph{Journal of Neuroscience}. Theoretical basis for smooth, jerk-minimizing movement.
\item[Flash \& Hogan (1985).] The coordination of arm movements: an experimentally
  confirmed mathematical model. \emph{Journal of Neuroscience}. Source of the
  fifth-order minimum-jerk trajectory and bell-shaped velocity profile.
\item[Xia \& Frey-Law (2008).] A theoretical approach for modeling peripheral
  muscle fatigue and recovery. \emph{Journal of Biomechanics}. Inspiration for a
  saturating fatigue state.
\item[Frey-Law \& Avin (2010).] Endurance time is joint-specific. \emph{Ergonomics}.
  Justifies variability in fatigue resistance.
\item[Frey-Law, Looft \& Heitsman (2012).] A three-compartment muscle fatigue model.
  \emph{Journal of Biomechanics}. Reference for fatigue accumulation/plateau behavior.
\item[Potvin (1997).] Effects of muscle kinematics on surface EMG amplitude and
  frequency during fatiguing dynamic contractions. \emph{J.\ Applied Physiology}.
  Dynamic fatigue signatures during elbow flexion--extension.
\item[Savin et al.\ (2021).] Movement variability patterns during a repetitive
  pointing task until exhaustion. \emph{Applied Ergonomics}. Basis for fatigue-related
  endpoint drift / movement variability.
\item[Yang et al.\ (2018).] Changes in movement variability and task performance
  during a fatiguing repetitive pointing task. \emph{Journal of Biomechanics}.
  Additional support for fatigue-related movement-variability changes.
\item[Attarieh et al.\ (2025).] Comparison of shoulder positions in elbow-flexion
  resistance training (preacher vs.\ Bayesian cable curls). \emph{European Journal
  of Sport Science}. Source for the controlled-curl elbow ROM ($10\dg$--$140\dg$)
  used to inform the chosen $20\dg$--$120\dg$ range.
\item[Sheikhhoseini et al.\ (2025).] Biomechanical coordination and variability after
  repetitive movement fatigue. \emph{Scientific Reports}. Post-fatigue movement-variability changes.
\end{description}

\end{document}
